%% kapa: exam=1 round=29 year=2018 period=1 form=A subject=civil_law subject_ko=민법 seq=32 number=32
%% source: Q-Net 감정평가사 기출문제 (cst00309) artlSeq=5202808 file=extracted/1008/5202808/2018년 제29회 감정평가사 1차시험 기출문제/2018년 제29회 감정평가사 1차시험 기출문제/1교시 A형.pdf
%% conversion: kordoc 4.24.51 -> tools/gen1.py (auto); review=proofread-sonnet; flags=none
\begin{question}{32}
점유에 관한 설명으로 옳지 않은 것은? (다툼이 있으면 판례에 따름)
\begin{choices}
  \item 점유자는 소유의 의사로 선의, 평온 및 공연하게 점유한 것으로 추정된다.
  \item 승계취득자가 전점유자의 점유를 아울러 주장하는 경우에는 그 점유의 하자도 승계한다.
  \item 임치관계로 타인으로 하여금 물건을 점유하게 한 자는 간접으로 점유권이 있다.
  \item 선의의 점유자라도 본권에 관한 소에 패소한 때에는 그 판결이 확정된 때로부터 악의의 점유자로 본다.
  \item 선의의 점유자는 비록 법률상 원인 없이 타인의 건물을 점유ㆍ사용하더라도 그로 인한 이득을 반환할 의무가 없다.
\end{choices}
\end{question}
\answer{4}
